%% ma: L12-B07-0034
%% lop: 12
%% chuong: 2
%% bai: 7
%% dang: 12.7.10
%% loai: TN4
%% muc: TH
%% dap_an: "A"
%% nguon: "CĐ2.6.pdf giáo viên gửi 27/10/2026 (M07, Lê Minh Khôi – Nguyễn Anh Khôi, Chương II Chủ đề 6 Ứng dụng vào các bài toán thực tế), Câu 29"
%% kien_thuc: "Bài 7 (tọa độ điểm khi gắn hệ trục)"
%% trang_thai: chinh-thuc
%% ghi_chu: "Giáo viên đã duyệt 27/10/2026. Đề gốc dùng chung giả thiết và hình cho Câu 29–30; đã chép vào từng câu"
\begin{ex}
Trong không gian $Oxyz$ (mỗi đơn vị dài trên các trục tương ứng $1$ mét, mặt phẳng ngang là mặt đất $(Oxy)$), cho một tháp chơi dành cho trẻ em leo trèo và trượt. Tháp hai tầng, tầng 1 là hình lập phương cạnh $2$ m, tầng hai có hình hộp chữ nhật cao $1$ m (cùng đáy với tầng 1) và một mái hình chóp tứ giác đều cao $1$ m. Từ đỉnh $S$ của tháp, một sợi dây được cố định tại điểm $P(1;7;0)$ để hỗ trợ leo trèo. Một bề mặt trượt hình chữ nhật được gắn vào mép tháp tại độ cao $2$ m sao cho mặt trượt song song với dây hỗ trợ leo và chạm đất (xem hình vẽ). Hệ trục tọa độ được chọn sao cho gốc $O(0;0;0)$ trùng với một góc của đáy tháp và đáy tháp nằm trong góc phần tư thứ nhất của mặt phẳng $(Oxy)$.
\begin{center}
\begin{tikzpicture}[scale=.85,x={(-.5cm,-.42cm)},y={(1cm,0cm)},z={(0cm,1cm)}]
  \draw[truc] (0,0,0)--(4,0,0) node[left]{$x$};
  \draw[truc] (0,2,0)--(0,8.5,0) node[right]{$y$};
  \draw[truc] (0,0,4)--(0,0,5.5) node[left]{$z$};
  \path (2,0,0) coordinate (A) (2,2,0) coordinate (B) (2,5,0) coordinate (C) (2,0,2) coordinate (E) (2,2,2) coordinate (F) (2,0,3) coordinate (H) (2,2,3) coordinate (G)
        (0,2,2) coordinate (F0) (0,5,0) coordinate (C0) (1,1,4) coordinate (S) (1,7,0) coordinate (P);
  \fill[gray!45] (F0)--(F)--(C)--(C0)--cycle; \draw (F0)--(F)--(C)--(C0);
  \draw (A)--(B)--(0,2,0) (A)--(H)--(G)--(B) (E)--(F) (G)--(0,2,3)--(0,2,0) (H)--(0,0,3)--(0,2,3) (S)--(H) (S)--(G) (S)--(0,2,3);
  \draw[khuat] (0,0,0)--(A) (0,0,0)--(0,2,0) (0,0,0)--(0,0,3) (S)--(0,0,3);
  \draw[red] (S)--(P); \fill (P) circle (2pt) node[right]{$P$};
  \draw[<->] (2,-.6,2)--(2,-.6,3) node[midway,left]{$1$ m};
  \path (S) node[above]{$S$} (A) node[below]{$A$} (B) node[below]{$B$} (C) node[below]{$C$} (E) node[below right]{$E$} (F) node[below right]{$F$} (H) node[above right]{$H$} (G) node[above right]{$G$} (0,0,0) node[above left]{$O$}
        ($(A)!.5!(B)$) node[below]{$2$m};
\end{tikzpicture}
\end{center}
Tọa độ của điểm $S$ là
\choice{\True $S(1;1;4)$}{$S(1;1;3)$}{$S(2;2;4)$}{$S(0;0;4)$}
\loigiai{Đáy tháp là hình vuông $[0;2]\times[0;2]$, tâm $(1;1)$; độ cao đỉnh $2+1+1=4$. Vậy $S(1;1;4)$.}
\end{ex}
